\documentclass{article}
\usepackage{graphicx} % Required for inserting images
\usepackage[utf8]{inputenc}
\usepackage[T2A]{fontenc}

\title{ML Project}
\date{December 2024}

\begin{document}


\section{Вопросы по типам признаков}
1. Категориальные признаки не могут сравниваться, в отличии от порядковых. То есть если мы возьмем один порядковый признак из одной группы и другой из той же группы, то мы точно сможем сравнить их. Чего мы не сможем сделать в случае категориальных \\\\
2. Распределение:
\begin{itemize}
    \item Категориальные признаки: is\_holiday, month
    \item Порядковые: weather\_condition
    \item Количественные: taxi\_price
\end{itemize}
weather\_condition я отнес к порядковым, т.к. погоду оценивали от 1 до 5, и мы точно можем сказать какая погода лучше а какая хуже \\
is\_holiday я отнес к категориальным, т.к. либо праздник либо нет, то есть признак бинарный, а бинарные почти всегда категориальные + мы не можем сравнить состояния "праздник" и "не праздник" \\
taxi\_price точно числовой признак, можем сравнить какая цена больше, какая меньше, можем спокойно вычитать и складывать, и выражен числом \\
month - категориальный, тк мы не можем сравнить два месяца, мы не поймем какой лучше, а какой хуже




\section{Определяем задачу машинного обучения}

\begin{enumerate}
    \item Объект: какой-то определенный самокат \\ 
    Целевая переменная: подозрительное ли поведение \\ 
    Тип: бинарная классификация (обучение с учителем) \\
    \item Объект: какой-то сотрудник из компании \\ 
    Целевая переменная: время в которое надо выходить из дома \\ 
    Тип: Регрессия (обучение с учителем)
\end{enumerate}

\section{Общие вопросы}
\begin{enumerate}
    \item Функция потерь - операция которая дает нам понять насколько то что выдает модель отличается от того что реально есть (мера между предсказанием и таргетом для 1 объекта), а Функционал ошибки использует эту функцию для обобщения и нахождения среднего по всем данным.
    \item Обучение модели это процесс когда мы имеем датасет, этот датасет мы кормим модели, а модель пытается выдать нам желаемый результат. Мы способны менять параметры этой модели и пытаемся поменять их таким образом чтобы модель ошибалась меньше всего.
\end{enumerate}


\end{document}

